%==============================================================================
\chapter{الوصف الوظيفي للحاجيات}
%==============================================================================

\section{إدارة الكيانات (CRUD)}
يجب أن يضمن النظام إنشاء جميع الكيانات المهنية واستعراضها وتعديلها وحذفها:
مربّو النحل والمزارع والمواقع والخلايا والأجسام/العاسلات والإطارات والأعوان
والزيارات والتقارير، مع احترام قواعد التسيير المذكورة في الفصل 3.

\section{تخطيط الزيارات ومتابعتها}
\begin{itemize}
  \item الزيارات \textbf{منتظمة أو غير منتظمة}، لأسباب متنوّعة: تفتيش صحي،
        تعمير بطرد جديد، تغيير الملكة، تزويد بالغذاء، وصف دواء، تقسيم الطرد،
        إضافة إطار و/أو طابق تمديد، جمع إطارات العسل، تقييم الإنتاج، تقييم عدد
        الطرد وحجمه.
  \item يتّبع كل عون \textbf{برنامج زيارة مُصادَقًا عليه} من العون المشرف على الخلية المعنية.
\end{itemize}

\subsection{تقرير الزيارة}
تُنتِج كل زيارة تقريرًا يحتوي على: \textbf{التاريخ، الساعة، المدة، السبب،
الملاحظات، الإجراءات المقرّرة، الإجراءات المنجَزة، التوصيات}، وتقييمًا نوعيًا
لعدد الطرد وحجمه وحالته الصحية ومستوى إنتاجيته على \textbf{سلّم من 1 إلى 3}.

\section{لوحات القيادة}

\subsection{الرزنامة / مصفوفة الأعوان $\times$ الخلايا}
مصفوفة تقاطع \textbf{الأعوان والخلايا}، مُقسّمة حسب الشهر والأسبوع والموسم
والسنة، مع \textbf{مرشّحات} وتجميعات حسب الموقع والعون والعون المسؤول. تُظهِر
النقرة على خانة تمثّل زيارة مبرمَجة \textbf{نافذة منبثقة} تُلخّص ما هو مقرّر
(التاريخ، الساعة، السبب، الإجراءات)، وإن كانت الزيارة قد نُفّذت، ما تمّ إنجازه
مع التاريخ والساعة الفعليين.

\subsection{لوحة قيادة الإنتاج (مربّي النحل / المالك)}
تُبرِز أزواج \textbf{الموقع $\times$ الخلية} التي تجاوز إنتاجها (الوزن)
\textbf{عتبة قابلة للضبط}، ما يشير إما إلى جمع عسل وشيك، أو إلى عملية تمديد و/أو
تقسيم طرد لتشجيع توسّع الطائفة.

\subsection{لوحة قيادة التنبيهات (المسؤول)}
تُشير إلى أزواج \textbf{الموقع $\times$ الخلية} التي ينخفض إنتاجها أو عددها بشكل
غير معتاد بما يتجاوز عتبة قابلة للضبط، مُطلِقةً \textbf{تنبيهًا} وتفتيشًا وشيكًا.

\subsection{لوحة القيادة الخلاصية (المالك)}
رسوم ومنحنيات ومدرّجات تكرارية تمثّل:
\begin{itemize}
  \item الكمية المُنتَجة حسب المزرعة والموقع
  \item عدد التدخّلات حسب المزرعة والموقع والخلية
  \item \textbf{العائد على الاستثمار} (نسبة الكلفة / الإنتاج) للفصل في مصير كل موقع أو خلية (الإغلاق، النقل).
\end{itemize}

\section{مؤشّرات الأداء (الشقّ الآلي: التهيئة)}
القياسات \textbf{مؤرّخة ومحدّدة الموقع}، مُعبَّر عنها بوحدات \textbf{قابلة
للضبط} (تدويل أنظمة الوحدات). لا تتشارك قياسات المؤشّر الواحد بالضرورة الوحدة
نفسها، لأن المستشعرات غير متجانسة. المؤشّرات المعنية:
\begin{itemize}
  \item وزن الإطار والقاعدة وطابق التمديد
  \item الحرارة والرطوبة، داخل الخلية وخارجها
  \item عدد حركات دخول النحل وخروجه
  \item مستوى الضجيج / الطنين، الداخلي والخارجي
  \item مستوى الإضاءة، الداخلي والخارجي
  \item سرعة الرياح واتجاهها في الخارج
  \item حالة فتح الخلية (تدخّل، عاصفة، سرقة، هجوم حيواني، إلخ).
\end{itemize}

\subsection{التحليل التنبّئي والمراقبة المتّصلة}
تُظهِر حلول المراقبة المتّصلة (تربية النحل الدقيقة) مثل \textenglish{Arnia} أو
\textenglish{BroodMinder} أو \textenglish{BeeHero} أن الاستغلال المتقدّم
للمؤشّرات نفسها يتيح استباق أحداث الطائفة. يجب أن يبقى نموذج المعطيات منفتحًا على
هذه المعالجات، المقرّرة للشقّ الآلي.
\begin{itemize}
  \item الكشف الصوتي عن التطريد من 1 إلى 5 أيام مسبقًا بتحليل البصمة الصوتية للطائفة
  \item تأكيد وجود الملكة انطلاقًا من الطنين
  \item الكشف عن فتح الخلية أو صدمها أو انقلابها بمقياس التسارع (سرقة، سقوط، هجوم حيواني)، إضافةً إلى مؤشّر الفتح
  \item تحديد اتجاه الخلية ببوصلة إلكترونية
  \item الوزن المستمر بميزان لتتبّع تدفّق الرحيق ودينامية الإنتاج
  \item التحليل السحابي بالذكاء الاصطناعي للبصمات الصوتية لربط الصوت بنشاط الخلية.
\end{itemize}

\subsection{بروتوكول استيعاب القياسات}
يبقى الشقّ الآلي خارج نطاق التطوير، لكن واجهة الاستيعاب مُحدَّدة منذ الآن كي
يستبقها نموذج المعطيات والبنية.
\begin{itemize}
  \item \textbf{النقل}: تُرسل بوّابة ميدانية القياسات إلى الخادم، حسب الاتصال
        المتاح، إما عبر \textbf{REST على HTTPS} (\texttt{POST /api/mesures}) أو
        عبر \textbf{MQTT} (نشر على الموضوع \texttt{zumm/\{rucheId\}/\{indicateur\}}).
  \item \textbf{الصيغة}: \textbf{JSON}، قياس يحمل معرّف الخلية ونوع المؤشّر
        والقيمة والوحدة والطابع الزمني والموقع الجغرافي. الإرسال بالحزمة عبارة عن جدول قياسات.
  \item \textbf{التواتر}: قابل للضبط حسب المؤشّر (مثلًا الوزن كل 15 دقيقة،
        الحرارة كل 5 دقائق، الصوت مُعايَن باستمرار)، للتحكّم في الحجم.
  \item \textbf{المتانة}: استيعاب \textbf{لا تزامني} عبر طابور انتظار، مع
        \textbf{عدم التكرار (idempotence)} على المفتاح (الخلية، المؤشّر، الطابع
        الزمني) لإعادة التشغيل دون ازدواج. في المناطق النائية، تخزين محلي مؤقّت
        على البوّابة ثم مزامنة مؤجّلة.
  \item \textbf{الأمان}: بوّابة مُوثَّقة (مفتاح أو شهادة)، قناة TLS. رفض القيم
        الخارجة عن المجال الفيزيائي المعقول (حاجز أمان).
\end{itemize}
\begin{verbatim}
POST /api/mesures        Content-Type: application/json
{
  "rucheId": 42, "typeIndicateur": "TEMPERATURE_INTERNE",
  "valeur": 34.8, "unite": "degC",
  "horodatage": "2026-07-09T14:05:00Z",
  "latitude": 36.806, "longitude": 10.181
}
\end{verbatim}

\subsection{العتبات الافتراضية ومنطق التنبيه}
تُخزَّن العتبات في جدول \texttt{SEUIL} (قابلة للضبط، الفصل 5). القيم أدناه
\textbf{قيم افتراضية مرجعية}، مستندة إلى البروتوكول النحلي (الفصل 12).
\begin{xltabular}{\textwidth}{>{\bfseries}p{4.4cm} X p{4.4cm}}
\toprule
\textbf{المؤشّر} & \textbf{العتبة الافتراضية} & \textbf{التنبيه المُطلَق} \\
\midrule
\endhead
الحرارة الداخلية (الحضنة) & أقل من 32\,\textdegree م أو أكثر من 36\,\textdegree م & خطر حراري على الحضنة \\
الرطوبة الداخلية & أكثر من 80\,\% & فرط رطوبة، خطر صحي \\
الوزن: ربح العاسلة & أكثر من عتبة الإنتاج & جمع عسل وشيك \\
الوزن: انخفاض مفاجئ & أكثر من 1.5\,كغ في أقل من ساعة & تطريد أو سرقة أو سقوط أو هجوم \\
الإنتاج أو العدد & انخفاض نسبي أكثر من 20\,\% مقارنةً بالفترة السابقة & تفتيش صحي وشيك \\
نشاط الدخول/الخروج & $\approx$\,0 في يوم مواتٍ & طائفة ضعيفة أو ملكة غائبة \\
حالة الفتح & مفتوحة خارج زيارة مبرمَجة & اقتحام أو اضطراب \\
البصمة الصوتية & نمط تطريد مُكتشَف & إنذار مسبق بالتطريد (1 إلى 5 أيام) \\
\bottomrule
\end{xltabular}
\begin{encadre}
\textbf{قاعدة التقييم:} عند كل قياس، تُحوَّل القيمة إلى الوحدة المرجعية (صيغة
التحويل، §5.3)، ثم تُقارَن بالعتبة القابلة للضبط. للحدّ من \emph{التنبيهات
الكاذبة}، لا يُطلَق تنبيه إلا إذا \textbf{استمرّ التجاوز على $N$ قياسات متتالية}
(مانع الارتداد / التخلفية). يبقى كل تنبيه \textbf{مساعدةً على القرار} يؤكّدها
تقرير الزيارة: يحتفظ العون بالمصادقة النهائية.
\end{encadre}

\section{جدول تركيبي للوظائف}
\begin{xltabular}{\textwidth}{>{\bfseries}p{3.4cm} X C{2.2cm}}
\toprule
\textbf{الوظيفة} & \textbf{الوصف} & \textbf{الأولوية} \\
\midrule
\endhead
CRUD الكيانات       & إدارة كاملة للنموذج المهني                        & عالية \\
برمجة الزيارات      & التخطيط + المصادقة من المشرف                      & عالية \\
تقارير الزيارة      & إدخال منظّم للملاحظات والإجراءات                  & عالية \\
رزنامة الأعوان×الخلايا & مصفوفة قابلة للتصفية مع نافذة ملخّص             & عالية \\
تنبيهات الإنتاج     & كشف تجاوز العتبات                                 & عالية \\
التنبيهات الصحية    & كشف الانخفاضات غير الطبيعية في العدد/الإنتاج      & عالية \\
الخلاصة والعائد     & رسوم الكلفة/الإنتاج، مساعدة على القرار            & متوسّطة \\
مؤشّرات المستشعرات  & نموذج معطيات جاهز للأتمتة                         & متوسّطة \\
\midrule
\multicolumn{3}{r}{\textit{\color{mielfonce} متطلبات تفادي العيوب (مستخلصة من حدود الحلول الموجودة)}} \\
\midrule
النشر الذاتي        & تثبيت J2EE دون اشتراك أو خدمة خارجية إجبارية      & عالية \\
الوضع دون اتصال     & إدخال دون اتصال ومزامنة مؤجّلة                    & عالية \\
مصادقة بسيطة        & OAuth Google ومسار إدخال ميداني سريع             & عالية \\
نمطية قابلة للضبط   & تفعيل الوحدات حسب الملف عبر ConfigZumm.ini        & متوسّطة \\
قابلية نقل المعطيات & تصدير CSV وTXT، وخدمة ويب ونسخ احتياطي، دون احتكار & عالية \\
مساعدة سياقية       & واجهة بسيطة ومتّسقة، توجيه المستخدم               & متوسّطة \\
تكافؤ الويب والجوّال & الوظائف نفسها على جميع الوسائط                    & متوسّطة \\
التدويل             & ملف XML بالفرنسية والإنجليزية والعربية، قابل للتوسّع & عالية \\
انفتاح عتادي        & مستشعرات غير متجانسة ووحدات قابلة للضبط           & متوسّطة \\
استيعاب لا تزامني   & قياسات مؤرّخة مؤجّلة، شقّ يدوي مستقلّ             & متوسّطة \\
تحقّق بشري          & عتبات قابلة للضبط وتأكيد بتقرير الزيارة           & عالية \\
أمان التبادلات      & SSL وشهادات X.509، وصول موثَّق                    & عالية \\
مساعدة على القرار   & تنبيهات إرشادية، يحتفظ العون بالمصادقة النهائية   & متوسّطة \\
\bottomrule
\end{xltabular}

\section{وظائف إضافية مستوحاة من حلول السوق}
لتعزيز سهولة الاستعمال والقيمة العملية، يستلهم النظام أفضل ممارسات تطبيقات تربية
النحل المرجعية مثل \textenglish{HiveTracks}. توسّع هذه الوظائف الحاجيات المذكورة
دون أن تعوّضها.
\begin{xltabular}{\textwidth}{>{\bfseries}p{4.2cm} X C{2.1cm}}
\toprule
\textbf{الوظيفة} & \textbf{الوصف} & \textbf{الأولوية} \\
\midrule
\endhead
صور التفتيش & إرفاق صورة أو أكثر بكل تقرير زيارة لمتابعة بصرية للطائفة والإطارات. & عالية \\
إدخال صوتي مُساعَد & إملاء محضر الزيارة، والنسخ النصي يغذّي حقول التقرير تلقائيًا. & منخفضة \\
سياق الطقس & عرض طقس الأسبوع المحلي انطلاقًا من موقع الموقع لتنوير قرارات التدخّل. & متوسّطة \\
خريطة وأنصاف أقطار الرعي & تحديد مواقع الخلايا على خريطة ورسم أنصاف أقطار الرعي (مثلًا 1 و2 و3 كم، وحدات قابلة للضبط) لتصوّر المحيط المتاح للنحل. & متوسّطة \\
قائمة مهام وتذكيرات & إدارة قائمة مهام ورزنامة تذكيرات (التغذية، التفتيش، حالة الملكة) كي لا يفوت أي موعد. & عالية \\
متابعة الملكة & تسجيل تاريخ حالة الملكة (الوجود، العمر، الاستبدال، التعليم) عبر الزيارات. & عالية \\
الحصاد ورمز QR & تسجيل حصاد العسل وتوليد رمز QR لكل دفعة يعرض ملاحظات التذوّق والمصدر والصور، للتتبّع والتثمين. & متوسّطة \\
وصول متعدّد الوسائط & توفير وصول عبر الويب وتجربة ملائمة للأجهزة المحمولة للإدخال الميداني. & متوسّطة \\
\bottomrule
\end{xltabular}
\begin{encadre}
\textbf{الاتّساق مع المتطلبات:} تعيد أنصاف أقطار الرعي والطقس استعمال الموقع
الجغرافي للمواقع وآلية الوحدات القابلة للضبط. تعتمد التذكيرات على برنامج الزيارة
الموجود، وتُثري متابعة الملكة تقرير الزيارة المُعرَّف سلفًا.
\end{encadre}

\section{وظائف إدارة متقدّمة (استلهام من Apiary Book وBeePlus)}
يُبرز تطبيقا إدارة تربية النحل \textenglish{Apiary Book} و\textenglish{BeePlus}
وظائف تنظيم وتتبّع تُكمّل نطاق النظام.
\begin{xltabular}{\textwidth}{>{\bfseries}p{4.2cm} X C{2.1cm}}
\toprule
\textbf{الوظيفة} & \textbf{الوصف} & \textbf{الأولوية} \\
\midrule
\endhead
الوضع دون اتصال والمزامنة & إدخال المعطيات الميدانية دون اتصال ومزامنتها تلقائيًا عند عودة الشبكة. & عالية \\
متابعة العلاجات & حفظ تاريخ الأدوية والعلاجات الصحية المطبَّقة على كل خلية. & عالية \\
إدارة العتاد والجرد & تتبّع العتاد (أجسام، عاسلات، إطارات، خلايا) وتخصيصه للمواقع. & متوسّطة \\
المتابعة المالية & تسجيل الكلفة والإيرادات حسب المزرعة والموقع والخلية، دعمًا لحساب العائد على الاستثمار. & متوسّطة \\
تصدير المعطيات & تصدير السجلّات بصيغتي CSV وTXT لتحليل خارجي. & متوسّطة \\
كشف الأمراض & المساعدة على تحديد الأمراض انطلاقًا من ملاحظات الزيارة. & منخفضة \\
النسخ الاحتياطي & نسخ المعطيات واستعادتها في السحابة. & متوسّطة \\
\bottomrule
\end{xltabular}
\begin{encadre}
\textbf{الاتّساق مع المتطلبات:} تغذّي المتابعة المالية لوحة قيادة العائد على
الاستثمار، ويمدّد تصدير CSV واجهة خدمة الويب، وتُثري متابعة العلاجات سبب زيارة
«وصف دواء».
\end{encadre}

\section{الحدود المعروفة للحلول الموجودة والمتطلبات لتفاديها}
يُبرز تحليل تجارب المستخدمين حول حلول السوق (\textenglish{HiveTracks}،
\textenglish{Apiary Book}، \textenglish{BeePlus}) وأنظمة المراقبة المتّصلة
(\textenglish{Arnia}، \textenglish{BroodMinder}، \textenglish{BeeHero}) حدودًا
متكرّرة. يضع الجدول التالي كل حدّ في مقابل المتطلب المعتمَد لتفاديه في Zümm.
\begin{xltabular}{\textwidth}{>{\bfseries}p{5.2cm} X}
\toprule
\textbf{الحدّ المُلاحَظ} & \textbf{المتطلب المعتمَد لتفاديه} \\
\midrule
\endhead
اشتراك مدفوع وجدار دفع & نظام قابل للنشر الذاتي بـ J2EE، دون اشتراك إجباري؛ تبقى جميع وظائف الإدارة متاحة. \\
تشغيل محدود دون اتصال & وضع قوي دون اتصال مع مزامنة مؤجّلة؛ لا يتطلّب الإدخال الميداني اتصالًا دائمًا. \\
إنشاء حساب مفروض واحتكاك & مصادقة بسيطة عبر OAuth Google ومسار إدخال سريع؛ سهولة الاستعمال أولًا. \\
تحميل زائد بوظائف غير مفيدة & وحدات قابلة للتفعيل حسب الملف وعبر ملف الإعداد ConfigZumm.ini، بفضل البنية ضعيفة الاقتران. \\
تبعية وفقدان المعطيات & معطيات يتحكّم بها المستخدم، تصدير CSV وTXT، وخدمة ويب ونسخ احتياطي، دون احتكار مُلكي. \\
واجهة كثيفة ومنحنى تعلّم & واجهة بسيطة ومتّسقة وسهلة مع مساعدة سياقية، وفق متطلبات سهولة الاستعمال. \\
وظائف غير متكافئة حسب المنصّة & تكافؤ وظيفي بين الوصول عبر الويب والجوّال بفضل البنية متعدّدة الطبقات. \\
حلول بالإنجليزية فقط & تدويل عبر ملف XML، بالفرنسية والإنجليزية والعربية على الأقل، ومنفتح على لغات أخرى. \\
كلفة عتاد المستشعرات المرتفعة & انفتاح على مستشعرات غير متجانسة وزهيدة مع وحدات قابلة للضبط، دون احتكار عتادي. \\
الاستقلالية والاتصال في المناطق النائية & استيعاب لا تزامني لقياسات مؤرّخة؛ يعمل الشقّ اليدوي دون مستشعرات. \\
تنبيهات كاذبة وموثوقية المستشعرات & عتبات قابلة للضبط وتحقّق بشري عبر تقرير الزيارة؛ تبقى التنبيهات مساعدةً على القرار. \\
سرّية المعطيات وأمانها & تشفير SSL وشهادات X.509، وصول موثَّق. \\
الإفراط في الثقة بالأتمتة & تموضع كمساعدة على القرار؛ يحتفظ عون تربية النحل بالمصادقة النهائية. \\
\bottomrule
\end{xltabular}
